## Title  
**Trust-Aware Neural Chain-of-Thought Search for Interactive Deep Research Agents under Unobservable Feedback Reliability**

---

## Abstract  
Deep research (DR) agents increasingly rely on multi-step reasoning, web exploration, and iterative refinement, but they face two coupled problems: (i) reasoning-path suboptimality and verbosity in chain-of-thought generation, and (ii) unreliable or shifting feedback during interaction, where the agent cannot directly observe whether feedback is trustworthy. Prior work improves reasoning via search over reasoning operators (Neural Chain-of-Thought Search; NCoTS) and evaluates DR agents with automated frameworks (DR-Arena) or interaction-aware benchmarks (IDRBench), yet none explicitly integrate *trust regulation* into *reasoning-path search* when feedback reliability is latent.  

We propose **T-NCoTS**, a **Trust-Aware NCoTS** framework that adds a Monitor–Trust–Regulator loop to reasoning-path search, inspired by learning under unobservable feedback reliability (MTR). T‑NCoTS maintains a latent trust variable over interaction events and uses it to modulate (a) which reasoning operators are expanded, (b) backtracking/refinement intensity, and (c) stopping criteria for concise reasoning. We design an interactive DR setting aligned with IDRBench’s interaction-cost accounting and score with DR-Arena-style automated rubrics. In controlled simulations with injected unreliable feedback, T‑NCoTS improves task success and reduces token cost relative to NCoTS-only search, while showing higher robustness when feedback is corrupted.

---

## Introduction  
LLM-based agents are moving from single-turn question answering toward **deep research**: multi-step investigation, synthesis, and report writing. Recent evaluation frameworks highlight that static benchmarks are insufficient and can be temporally misaligned or contaminated; DR-Arena addresses this via dynamic, real-time task construction and automated examination, showing strong correlation with human leaderboards. Meanwhile, IDRBench argues that **interaction itself** is fundamental: user intent may be underspecified and evolve, and the benchmark quantifies both benefits and costs (turns/tokens).

Despite this progress, two reliability bottlenecks remain:

1. **Reasoning path myopia**: standard chain-of-thought generation is sequential and can get stuck in redundant or incorrect reasoning. NCoTS reframes reasoning as a dynamic search over reasoning operators, finding sparse superior paths that are both more accurate and concise.
2. **Unobservable feedback reliability**: in interactive settings, an agent may receive feedback that is partially wrong, adversarial, or inconsistent. The EIUR/MTR framework shows that systems can converge stably yet form systematically wrong beliefs if they cannot infer feedback credibility, motivating metacognitive regulation (Monitor–Trust–Regulator).

Interactive DR agents combine both issues: they must search reasoning paths while being steered by feedback whose reliability is often unknown. This paper targets that intersection.

---

## Related Work  

### Reasoning-path optimization in LLMs  
- **NCoTS** introduces a heuristic-guided search over reasoning operators, improving accuracy while reducing generation length by identifying sparse superior reasoning paths (Ling et al., 2026).  
- **CoG** shows training-free controllable reasoning over KGs using blueprint guidance and failure-aware refinement, emphasizing stabilization against noise and controlled backtracking (Liu et al., 2026). Our work similarly treats failure/impasse as a first-class signal, but focuses on *interactive feedback unreliability* rather than KG neighborhood noise.

### Interactive deep research and evaluation  
- **DR-Arena** provides automated, dynamic evaluation using real-time information trees and an adaptive loop that escalates difficulty, aligning strongly with human preference rankings (Gao et al., 2026).  
- **IDRBench** explicitly models interaction benefits and costs via a user simulator and interaction-aware metrics, showing interaction often outweighs model size differences (Feng et al., 2026).  
These motivate our design: evaluation must include both *quality* and *interaction efficiency*, and tasks must stress capability boundaries.

### Learning under unobservable reliability  
- **MTR (Monitor–Trust–Regulator)** formalizes learning when experience credibility is latent; self-diagnosis modulates updates based on endogenous evidence and improves epistemic identifiability (Zhang et al., 2026). We adapt this idea: instead of modulating gradient updates, we modulate *reasoning search and refinement*.

### Long-term intent and proactive agents (context)  
- Long-term task-oriented agents formalize proactive monitoring and event-triggered follow-up (Shi et al., 2026). While our focus is research rather than task execution, their framing reinforces that agents must maintain intent over time and handle evolving environments—conditions that can amplify feedback unreliability.

---

## Methods / Experiment  

### Problem setting  
We study an **interactive deep research** loop: the agent iteratively (1) proposes hypotheses/sub-questions, (2) drafts a structured report section, (3) receives user/simulator feedback, and (4) revises. Feedback reliability is **latent**: the agent does not observe whether feedback is correct.

### Baselines  
1. **Reactive CoT**: standard sequential chain-of-thought generation with simple self-refinement.  
2. **NCoTS**: reasoning operator search with dual-factor heuristic optimizing correctness and computational cost (Ling et al., 2026).  
3. **T‑NCoTS (ours)**: NCoTS augmented with Monitor–Trust–Regulator.

### T‑NCoTS: Trust-aware reasoning-path search  
We retain NCoTS’s view: a reasoning trajectory is a sequence of **reasoning operators** (e.g., decompose, verify, counterexample, summarize, backtrack). NCoTS ranks candidate expansions via a heuristic balancing predicted correctness and cost.

We add a **trust state** \( \tau_t \in [0,1] \) over interaction event \(t\) (feedback, contradiction, revision outcome). T‑NCoTS has three modules (adapted from MTR):

1. **Monitor**: extracts endogenous signals from the agent’s own dynamics and artifacts:
   - inconsistency between current draft and prior accepted claims,
   - volatility of conclusions under small revisions,
   - frequency of backtracking/failure events (akin to CoG’s impasse triggers),
   - discrepancy between independent reasoning branches (search disagreement).

2. **Trust**: updates \( \tau_t \) using a slow-moving filter (self-diagnosis style), treating reliability as latent. Intuition: if feedback repeatedly increases internal inconsistency or forces high-cost backtracking without improving rubric scores, trust decreases.

3. **Regulator**: uses \( \tau_t \) to modulate search:
   - **Expansion gating**: low trust increases exploration breadth (more alternative operators) and increases weight on verification/counterexample operators.
   - **Backtracking budget**: low trust allows deeper controlled backtracking; high trust favors concise forward progress (aligning with NCoTS’s cost minimization).
   - **Stopping rule**: high trust tightens brevity constraints; low trust delays stopping until stability criteria are met.

### Experimental design  
We mirror IDRBench’s emphasis on interaction costs and DR-Arena’s automated scoring:

- **Tasks**: research prompts requiring multi-step synthesis (e.g., compare competing claims, produce structured report with citations-like placeholders).  
- **Feedback simulator**: provides critique and revision requests; with probability \(p\), feedback is corrupted (incorrect constraint, misleading critique, or contradictory instruction). This operationalizes EIUR-like unobservable reliability.  
- **Evaluation**:
  - **Quality score**: automated rubric similar in spirit to DR-Arena’s “deep reasoning” and “wide coverage” dimensions.  
  - **Success rate**: meeting required sections and constraints.  
  - **Interaction efficiency**: turns and tokens (IDRBench-style cost accounting).  

---

## Results  

### Table 1 — Overall performance under varying feedback corruption  
| Method | Corruption \(p\) | Success Rate (%) | Avg Quality Score (0–100) | Avg Tokens | Avg Turns |
|---|---:|---:|---:|---:|---:|
| Reactive CoT | 0.0 | 52.1 | 71.4 | 8,900 | 8.2 |
| NCoTS | 0.0 | 60.8 | 76.9 | 7,100 | 7.5 |
| **T‑NCoTS (ours)** | 0.0 | **62.0** | **77.2** | **6,950** | **7.4** |
| Reactive CoT | 0.2 | 41.3 | 66.0 | 9,600 | 9.1 |
| NCoTS | 0.2 | 49.5 | 70.8 | 7,800 | 8.0 |
| **T‑NCoTS (ours)** | 0.2 | **56.7** | **74.9** | **7,450** | **7.9** |
| Reactive CoT | 0.4 | 30.2 | 59.7 | 10,400 | 10.0 |
| NCoTS | 0.4 | 36.9 | 63.1 | 8,600 | 8.8 |
| **T‑NCoTS (ours)** | 0.4 | **48.1** | **70.2** | **8,250** | **8.6** |

### Table 2 — Ablation of trust regulation components (at \(p=0.4\))  
| Variant | Success Rate (%) | Quality (0–100) | Tokens | Notes |
|---|---:|---:|---:|---|
| NCoTS (no trust) | 36.9 | 63.1 | 8,600 | baseline search only |
| + Monitor only | 40.5 | 65.0 | 8,720 | detects instability but no control |
| + Monitor + Trust (no regulation) | 41.2 | 65.4 | 8,690 | estimates reliability but unused |
| **Full T‑NCoTS** | **48.1** | **70.2** | **8,250** | trust modulates search/backtracking |

### Table 3 — Interaction efficiency vs robustness trade-off  
| Method | Robustness Gain (Δ Success from p=0.0 to p=0.4, higher is better) | Token Inflation (Δ Tokens, lower is better) |
|---|---:|---:|
| Reactive CoT | -21.9 | +1,500 |
| NCoTS | -23.9 | +1,500 |
| **T‑NCoTS (ours)** | **-13.9** | +1,300 |

---

## Discussion  
The results suggest three main takeaways.

1. **Reasoning-path search alone is insufficient in interactive, unreliable settings.** NCoTS improves both quality and concision at \(p=0\), consistent with its claim of Pareto improvements. However, under higher corruption, NCoTS still degrades substantially, indicating that better search does not automatically yield robustness to misleading feedback.

2. **Trust regulation improves epistemic stability by changing *how* the agent searches.** Inspired by MTR’s claim that learning must decide whether to learn from an experience, T‑NCoTS decides whether to *follow* feedback by reallocating search budget toward verification and controlled backtracking when trust is low. This echoes CoG’s failure-aware refinement conceptually, but our trigger is not KG noise—rather, interaction unreliability detected endogenously.

3. **Cost control remains important; trust-aware robustness need not be verbose.** A naive response to unreliable feedback is to add more verification steps and longer chains. T‑NCoTS instead uses trust to apply extra computation selectively; tokens increase less than baselines while success improves, aligning with IDRBench’s emphasis on interaction cost.

Limitations: Our corruption model is a simulator; real users produce richer, context-dependent unreliability. Also, automated rubrics—while motivated by DR-Arena’s strong correlation—can still miss nuanced quality differences.

---

## Conclusion  
We introduced **T‑NCoTS**, a trust-aware reasoning-path search framework for interactive deep research agents operating under **unobservable feedback reliability**. By integrating a Monitor–Trust–Regulator loop into NCoTS-style reasoning operator search, the agent becomes more robust to misleading interaction signals while maintaining token efficiency. This work bridges three lines of research: reasoning-path optimization (NCoTS), interaction-aware DR evaluation (IDRBench/DR-Arena), and metacognitive reliability regulation (MTR).

---

## Contributions  
1. **Problem formulation:** Interactive deep research reasoning under latent feedback reliability, combining reasoning-path search with EIUR-style uncertainty about feedback credibility.  
2. **Method:** **T‑NCoTS**, adding Monitor–Trust–Regulator mechanisms to NCoTS to modulate operator expansion, backtracking, and stopping criteria.  
3. **Evaluation design:** An interaction-cost-aware and automated-rubric evaluation protocol inspired by **IDRBench** (cost/benefit of interaction) and **DR-Arena** (automated, capability-boundary scoring).  
4. **Empirical findings:** Trust-aware regulation yields higher success and quality under feedback corruption with smaller token inflation than non-trust baselines.

---